\documentclass[11pt]{article}
\usepackage[margin=1in]{geometry}
\usepackage{booktabs}
\usepackage{graphicx}
\usepackage{amsmath}
\usepackage{hyperref}
\usepackage{natbib}
\usepackage{float}

\title{CB-Signal: Extracting a Monetary-Policy Tone Signal from\\
Central Bank Communications}
\author{Marouane Sakai}
\date{\today}

\begin{document}
\maketitle

\begin{abstract}
We build an end-to-end pipeline that scrapes European Central Bank (ECB),
Federal Reserve (Fed) and Bank of England (BoE) communications, scores
their hawkish/dovish tone with three independent methods (a curated
phrase dictionary, a pretrained financial-sentiment transformer, and a
lightweight fine-tuned sentence-embedding classifier), aggregates tone
into a weekly index, and backtests a resulting trading signal against
2Y/10Y government bond yields. On 3,968 real documents spanning
1997--2026, the baseline dictionary signal shows no robust full-sample
edge (Sharpe ratios from $-1.6$ to $+0.07$ across legs), but sub-period
analysis reveals the signal performs materially better during the
2022--2023 hiking cycle (Sharpe up to $1.8$) than in other regimes. On a
common window comparing all three sentiment tiers, the simple dictionary
outperforms both the pretrained and fine-tuned transformer tiers --- a
caution against assuming more sophisticated models automatically win on a
narrow, low-data task. We also document and fix a look-ahead issue in the
original regime-labelling procedure, and report where a walk-forward
refit changes the resulting labels. We report all of this plainly,
including the negative full-sample result, rather than presenting only
the sub-period in which the signal happens to work.
\end{abstract}

\section{Motivation}

Central bank communications are a primary channel through which monetary
policy expectations are formed. A large academic literature
\citep{apel2012information,hansen2016shocking,bennani2017home} argues that
the \emph{tone} of this language --- hawkish (tightening-oriented) versus
dovish (easing-oriented) --- carries information beyond the policy rate
decision itself. This project asks a narrower, falsifiable question: does a
systematic, text-derived tone signal, built with openly documented methods
and tested on real market data, produce a backtestable edge on government
bond yields once realistic transaction costs are applied?

\section{Data}

\subsection{Communications corpus}

We scrape three sources directly from the issuing institutions' own
websites (no third-party aggregator):

\begin{itemize}
  \item \textbf{ECB}: 1,871 speeches (1997--2026), discovered via the
    ECB's AddSearch-powered site search (the previously published bulk CSV
    feed was discontinued following a site redesign) and parsed from each
    speech's canonical HTML page.
  \item \textbf{Fed}: 1,649 documents --- 167 FOMC statements, 150 FOMC
    minutes, and 1,332 speeches (2006--2026) --- discovered via the Fed's
    historical and current meeting calendars and speech archives.
  \item \textbf{BoE}: 448 documents --- 371 speeches and 77 Monetary
    Policy Summary/minutes releases (2015--2026) --- discovered via BoE's
    sitemap feed, after the browsable listing pages moved to a
    JavaScript-rendered search widget.
\end{itemize}

\noindent Table~\ref{tab:corpus} summarises coverage. All ingestion code
retries transiently-failing requests, caches raw HTML on disk, and is
designed to degrade gracefully (skip, don't crash) when an individual page
is unreachable.

\begin{table}[H]
\centering
\begin{tabular}{lrrr}
\toprule
Source & Documents & Event types & Date range \\
\midrule
ECB & 1,871 & speech & 1997--2026 \\
Fed & 1,649 & statement, minutes, speech & 2006--2026 \\
BoE & 448 & speech, minutes & 2015--2026 \\
\midrule
Total & 3,968 & & 1997--2026 \\
\bottomrule
\end{tabular}
\caption{Real communications corpus used in this study.}
\label{tab:corpus}
\end{table}

\subsection{Market data}

Weekly 2-year and 10-year US Treasury yields (FRED series \texttt{DGS2},
\texttt{DGS10}), German and UK long-term government yields, and the VIX
are pulled directly from the FRED API (39,665 observations, 1956--2026).

\section{Methodology}

\subsection{Sentiment scoring tiers}

We score every document with three independent methods so that the
dictionary baseline's results can be cross-checked rather than taken on
faith:

\begin{enumerate}
  \item \textbf{Dictionary} (baseline): a curated hawkish/dovish phrase
    lexicon drawn from \citet{apel2012information},
    \citet{hansen2016shocking} and \citet{bennani2017home}, plus a
    Loughran--McDonald \citep{loughran2011liability} uncertainty subset,
    matched with negation handling over a 3-token window. Net tone is
    $(\text{hawkish hits} - \text{dovish hits}) / \text{total hits}$.
  \item \textbf{Pretrained transformer}: ProsusAI/FinBERT
    \citep{araci2019finbert}, a BERT model fine-tuned on financial news
    sentiment (positive/negative/neutral), with positive/negative remapped
    to dovish/hawkish. Long documents are chunked into overlapping
    512-token windows and the resulting per-chunk probabilities averaged.
    (The FOMC-specific \citet{shah2023trillion} RoBERTa model would be a
    closer match to this task, but its HuggingFace repository is
    access-gated; FinBERT is used as the publicly reproducible
    alternative.)
  \item \textbf{Fine-tuned tier}: sentence embeddings from
    \texttt{all-MiniLM-L6-v2} \citep{reimers2019sentencebert} with a
    multinomial logistic-regression head, trained on a small
    (482-example) hawkish/dovish/neutral set. \emph{This set is not an
    independently human-annotated research dataset} --- it is built by
    dropping each phrase from the dictionary lexicon into several carrier
    sentences, plus a hand-written list of genuinely neutral central-bank
    boilerplate sentences for the neutral class. We report it as a
    lightweight, transparently-constructed classifier, not as a validated
    gold-standard benchmark.
\end{enumerate}

\subsection{Signal construction}

Document-level tone is averaged per source per week, then averaged across
sources into a single weekly hawkish index. An exponential moving average
(12-week half-life) serves as the baseline; the surprise (index $-$ EMA) is
expanding-window z-scored (so, by construction, only information available
up to that week is used --- no look-ahead) and clipped to $[-3, 3]$. The
tradable signal is a 4-week rolling average of a 50/50 blend of the
z-scored level and the z-scored surprise, applied with full weight to the
2-year leg and 60\% weight to the 10-year leg.

\subsection{Regime labelling: a look-ahead issue, found and fixed}

A 3-state Gaussian HMM on the level and week-over-week change of the
hawkish index labels each week dovish/neutral/hawkish for descriptive and
diagnostic purposes. \emph{This regime label does not feed the tradable
signal above} --- the signal is built only from the already point-in-time
EMA/surprise construction. The original implementation fit the HMM once on
the entire sample, meaning early-sample regime labels were informed by
the model's fit to \emph{later} data: a look-ahead issue for a
descriptive label, even though it has no effect on realised backtest P\&L.
We implemented a walk-forward version that refits the HMM yearly on an
expanding window and labels only the following year out-of-sample. The
walk-forward and in-sample labels disagree on 45\% of weeks, showing the
original full-sample fit was not a robust description of the regime
structure, and this out-of-sample version should be preferred whenever the
regime label itself is used for anything beyond the illustrative plot
included here.

\subsection{BERTopic thematic decomposition}

As a complementary descriptive layer, we fit BERTopic
\citep{grootendorst2022bertopic} on sentence embeddings of the full corpus
to recover the dominant themes discussed across communications (inflation,
labour markets, financial stability, etc.) and their prevalence over time.
This is reported alongside the tone signal as context, not as an input to
it.

\subsection{Backtest}

Weekly yield changes are converted to approximate bond returns via a
constant modified-duration approximation ($D_{2Y}=1.9$, $D_{10Y}=8.5$).
Positions are $-\mathrm{sign/magnitude}$ of the (lagged, $T+1$) signal,
clipped to $[-1,+1]$, with 20bps round-trip transaction costs charged on
turnover.

\section{Results}

\subsection{Full-sample tearsheet}

Table~\ref{tab:main} reports full-sample results for the dictionary tier
(weekly, 1997--2026, 20bps round-trip TC, $T{+}1$ execution). FinBERT and
fine-tuned-tier results on the full corpus are not reported here --- at
this project's measured throughput on CPU (about 8 seconds per document
for FinBERT), scoring the full 3,968-document corpus would take roughly
9 hours; Section~\ref{sec:tiers} instead compares all three tiers on a
smaller, tractable window.

\begin{table}[H]
\centering
\begin{tabular}{lrrrrr}
\toprule
Strategy & Ann.\ return & Ann.\ vol & Sharpe & Max DD & Turnover/yr \\
\midrule
Short 10Y   & $+0.07\%$ & $0.92\%$ & $0.07$  & $-6.1\%$  & 2.2 \\
50/50       & $-0.24\%$ & $0.60\%$ & $-0.40$ & $-10.0\%$ & 2.9 \\
Short 2Y    & $-0.55\%$ & $0.34\%$ & $-1.61$ & $-15.2\%$ & 3.6 \\
\bottomrule
\end{tabular}
\caption{Dictionary-tier full-sample results on real data.}
\label{tab:main}
\end{table}

None of the three legs clears a Sharpe ratio of 1 over the full sample;
the short-2Y leg is clearly negative. This is a materially weaker result
than the same pipeline run on the project's synthetic data generator
(Sharpe up to 1.75), which is expected --- the synthetic generator ties
tone to returns by construction, real markets do not oblige.

\subsection{Sub-period stability}

Table~\ref{tab:subperiods} breaks the same backtest into five named
macro-regimes.

\begin{table}[H]
\centering
\small
\begin{tabular}{llrrrr}
\toprule
Period & Strategy & Ann.\ return & Sharpe & Max DD & Hit rate \\
\midrule
GFC (2007--2009)         & Short 2Y  & $-0.92\%$ & $-2.80$ & $-2.7\%$  & 0.25 \\
                         & Short 10Y & $-0.51\%$ & $-0.65$ & $-1.8\%$  & 0.41 \\
                         & 50/50     & $-0.71\%$ & $-1.35$ & $-2.2\%$  & 0.39 \\
ZIRP (2009--2015)        & Short 2Y  & $-0.75\%$ & $-4.16$ & $-5.1\%$  & 0.21 \\
                         & Short 10Y & $-0.25\%$ & $-0.35$ & $-2.3\%$  & 0.45 \\
                         & 50/50     & $-0.50\%$ & $-1.18$ & $-3.5\%$  & 0.37 \\
COVID (2020--2021)       & Short 2Y  & $-0.44\%$ & $-2.90$ & $-1.0\%$  & 0.18 \\
                         & Short 10Y & $-0.20\%$ & $-0.37$ & $-1.2\%$  & 0.40 \\
                         & 50/50     & $-0.32\%$ & $-1.01$ & $-1.1\%$  & 0.38 \\
\textbf{Hiking (2022--2023)} & \textbf{Short 2Y}  & $\mathbf{+0.80\%}$ & $\mathbf{1.55}$ & $-0.4\%$ & 0.49 \\
                         & \textbf{Short 10Y} & $\mathbf{+2.37\%}$ & $\mathbf{1.79}$ & $-0.9\%$  & 0.54 \\
                         & \textbf{50/50}     & $\mathbf{+1.58\%}$ & $\mathbf{1.83}$ & $-0.6\%$  & 0.57 \\
Plateau (2024+)          & Short 2Y  & $-0.85\%$ & $-2.58$ & $-2.5\%$  & 0.27 \\
                         & Short 10Y & $-0.31\%$ & $-0.42$ & $-1.9\%$  & 0.46 \\
                         & 50/50     & $-0.58\%$ & $-1.12$ & $-2.1\%$  & 0.39 \\
\bottomrule
\end{tabular}
\caption{Dictionary-tier sub-period stability. The signal's only
consistently positive period is the 2022--2023 hiking cycle.}
\label{tab:subperiods}
\end{table}

This is one of the paper's most interesting findings: the full-sample result
is negative primarily because of ZIRP, GFC, COVID and the current plateau
period, while the 2022--2023 hiking cycle --- arguably the regime where
central bank tone genuinely diverged most sharply from where it had been
anchored for over a decade --- is the one period where the signal
performs well across every leg. This is consistent with a hypothesis that
textual tone carries more information when the policy stance is actually
\emph{changing} quickly, and less when it is anchored near a bound (ZIRP)
or dominated by an exogenous shock (COVID). We present this as a
hypothesis the sub-period table motivates, not as a validated finding ---
one hiking cycle is one data point.

\subsection{Cross-tier comparison}
\label{sec:tiers}

Table~\ref{tab:tiers} compares all three sentiment tiers on a common,
tractable sample: a contiguous 2021--2023 window (688 documents --- a
random subsample across the full history would break the weekly
aggregation continuity the signal construction relies on), chosen to
include the 2022--2023 hiking cycle identified above plus a year of
run-in for the EMA/z-score to stabilise.

\begin{table}[H]
\centering
\begin{tabular}{lrrr}
\toprule
Tier & Short 10Y Sharpe & 50/50 Sharpe & Short 2Y Sharpe \\
\midrule
Dictionary  & $\mathbf{1.21}$  & $\mathbf{0.93}$  & $0.02$  \\
FinBERT     & $0.03$  & $-0.34$ & $-1.16$ \\
Fine-tuned  & $-0.63$ & $-1.36$ & $-2.89$ \\
\bottomrule
\end{tabular}
\caption{Sentiment-tier comparison, common 2021--2023 window.}
\label{tab:tiers}
\end{table}

On this window the purpose-built dictionary lexicon outperforms both
learned tiers, not the other way around. We read this as plausible rather
than surprising: FinBERT is trained on general financial-news sentiment,
not central-bank language specifically, and the fine-tuned tier's training
set is small, templated and author-curated (Section~3.1) rather than a
large corpus of naturally-occurring, independently-labelled central-bank
text. The result is a useful caution against assuming a larger pretrained
model automatically outperforms a smaller, domain-specific lexicon on a
narrow task with little in-domain labelled data.

\section{Limitations}

\begin{itemize}
  \item \textbf{No robust full-sample edge.} The headline result is
    negative-to-flat. We report it because it is true, not because it is
    the more publishable finding.
  \item \textbf{Small fine-tuned set.} The 482-example fine-tuned tier is
    author-curated and templated, not independently annotated; its train
    accuracy (95\%) reflects fitting a near-separable templated set, not
    held-out generalisation to real prose.
  \item \textbf{One hiking cycle.} The sub-period finding rests on a
    single historical episode (2022--2023); it is suggestive, not
    statistically established.
  \item \textbf{Duration proxy.} Returns use a constant modified-duration
    approximation rather than actual bond/future contract specifications;
    real deployment, particularly on the 2-year leg, would likely use
    short-term interest rate futures (e.g.\ SR3) rather than cash bonds,
    where round-trip costs are closer to 1bp than the 20bps assumed here.
  \item \textbf{Single aggregate rates proxy per currency area.} We do not
    model the full yield curve or cross-currency spillovers between the
    three central banks' communications.
\end{itemize}

\section{Conclusion}

A reproducible, real-data pipeline from central bank website to backtest
tearsheet does not, on its own, produce a profitable signal with a simple
dictionary-tone baseline. It does surface a specific, falsifiable
hypothesis --- tone-based signals may work better in regime-change periods
than in anchored or shock-dominated ones --- worth testing against a
larger sample of historical hiking/cutting cycles than the single episode
available in this dataset. Fully reproducing every number in this paper
takes under two hours end to end:
\texttt{python -m cb\_signal.ingest \&\& python -m cb\_signal.run\_all
--skip-synth --subperiods --walk-forward}.

\bibliographystyle{plainnat}
\bibliography{references}

\end{document}
